\section{Materiais e Métodos}

\subsection{Visão geral}

O pipeline possui cinco etapas: (i) descoberta de instâncias do SAPL a partir da lista de municípios do IBGE \cite{ibge_localidades_api}; (ii) extração de Projetos de Lei diretamente das APIs públicas; (iii) tradução de ementas em ações operacionais; (iv) indexação vetorial para busca semântica; e (v) associação exploratória com indicadores municipais e agrupamento de políticas.

\subsection{Descoberta e extração}

A descoberta combina heurísticas de URL com validação por evidência pública. Para cada município, o sistema gera candidatos como \texttt{sapl.\{slug\}.\{uf\}.leg.br}, testa rotas típicas do SAPL e valida a instância quando encontra marcadores compatíveis. Depois disso, a extração deriva a base correta da API, consulta os tipos de matéria e percorre a paginação de todos os tipos cujo rótulo contém ``projeto de lei''.

Na execução congelada em 12/11/2025, foram encontradas 1.259 instâncias do SAPL e 322 delas permitiram extração válida com ao menos um PL, totalizando 220.065 registros. Considerando os 5.570 municípios brasileiros, isso corresponde a 22,60\% de cobertura na descoberta e 5,78\% na extração automatizada. O gap entre essas etapas reflete heterogeneidade técnica entre portais, mudanças de rota, autenticação e diferenças locais de configuração.

O dataset final cobre 322 municípios em 19 UFs. Todos os registros possuem \textit{ementa}, \textit{ação}, \textit{link\_publico} e embedding preenchidos. A base não apresentou duplicatas por \textit{link\_publico} nem por \texttt{(sapl\_base, materia\_id)}. Por outro lado, o filtro textual de tipos ainda recupera cerca de 3,54\% de itens que não são PLs estritos, como anteprojetos, emendas, substitutivos e vetos.

\subsection{Tradução de ementas em ações}

Para reduzir ruído jurídico e tornar as proposições mais comparáveis, cada ementa é convertida em uma ação no formato ``verbo no infinitivo + objeto + complementos essenciais''. O conjunto de treinamento foi composto por 1.000 ementas reais distintas, textualizadas com apoio de GPT-5.1-Thinking \cite{openai_gpt51_model,openai_gpt51_system_card} e revisadas manualmente nos casos de perda de sentido.

Como modelo base, utilizou-se o PTT5-v2 \cite{piau_ptt5v2_2024}, uma variante do T5 \cite{raffel_t5_2020} adaptada ao português. O ajuste fino foi realizado com QLoRA em 4 bits \cite{dettmers_qlora_2023}, com divisão 90/10 entre treino e validação, 30 épocas, taxa de aprendizado de $3\times 10^{-4}$ e \textit{batch size} 16. A avaliação automática por BERTScore \cite{zhang_bertscore_2019} atingiu 84\% no conjunto de validação.

No acervo completo, a tradução reduziu o comprimento médio das descrições de 152,86 para 90,01 caracteres e de 24,18 para 13,61 palavras. Isso corresponde a 58,9\% dos caracteres e 56,3\% das palavras das ementas originais, favorecendo busca e agrupamento. O principal \textit{baseline} considerado até aqui foi indexar a ementa bruta.

\subsection{Busca, indicadores e agrupamento}

As ações são indexadas com embeddings do modelo Multilingual E5 \cite{wang_e5_2024} e armazenadas no Qdrant \cite{qdrant_overview,qdrant_collections}. As consultas usam o prefixo \texttt{query:}, conforme a recomendação do modelo. A avaliação da recuperação ainda é exploratória: há inspeção manual das primeiras posições, mas ainda não há benchmark anotado nem métricas formais como Precision@K ou nDCG.

Para aproximar o sistema de desfechos observáveis, utilizamos taxa de homicídios por 100 mil habitantes e taxa de matrículas em ensino regular por 100 mil habitantes. Os efeitos são calculados comparando o indicador na data mais próxima ao PL com janelas posteriores de 6, 12, 18, 24, 30 e 36 meses. Trata-se de análise descritiva, sem pretensão causal.

PLs semanticamente próximos são agrupados por similaridade de Jaccard sobre o texto da ação, com limiar 0,75. Para cada grupo, definimos o escore heurístico:

\[
Q = \frac{n_{\text{positivos}}}{n} \times \frac{n}{n+1}
\]

em que $n$ é o número de municípios representados e $n_{\text{positivos}}$ o número de efeitos favoráveis ao objetivo do indicador. O termo final penaliza grupos muito pequenos e ajuda a priorizar políticas mais recorrentes.

\subsection{Implementação}

A plataforma foi implementada em React/Next.js e integra interface web, serviço de embeddings, banco vetorial e dados de indicadores \cite{vercel_nextjs}. O código de coleta, treinamento, indexação e interface permanece disponível em \url{https://github.com/thiagoambiel/SonarMunicipal}, juntamente com o prompt do tradutor e a lista de instâncias SAPL utilizadas nesta versão.
